\documentclass{article}
\usepackage{tikz}
\usetikzlibrary{angles, quotes}
\usepackage{amssymb}
\usepackage{pgfplots}
\pgfplotsset{width=10cm,compat=1.9}
\usepgfplotslibrary{statistics}
\usepackage[left=2.54cm,right=2.54cm,top=2.54cm,bottom=2.54cm]{geometry}
\usepackage{tasks} %Paquete necesario para la producción de items horizontales para algunos ejercicios de tarea empleando el comando /begin{multicols}}{#}.../end{multicols} para ayudarnos a reducir las páginas
\usepackage{pdflscape} %comando que permite cambiar a orientación horizontal las hojas%
\usepackage{soul} %Paquete para permitir la compilación de acentos usando el comando {\hl{}}}\\
\sethlcolor{yellow(munsell)} %comando para definir el color para el subrayado usando el comando \hl
\usepackage[utf8]{inputenc}
\usepackage[spanish]{babel}
\usepackage{icomma}
\usepackage{ragged2e}		
\usepackage{siunitx}
\usepackage{url}
\usepackage[colorlinks=true, urlcolor=blue,  linkcolor=Cerulean, citecolor=green]{hyperref}
\usepackage{pdfpages}
\usepackage{blindtext} %Paquete que genera texto ficticio (dummy text) con el comando: \blindtext
\usepackage{cancel} %in the preamble gives you four different modes of striking through: \cancel{text to cancel} draws a diagonal line (slash) through its argument, \bcancel{text to cancel} uses the negative slope (a backslash), \xcancel{text to cancel} draws an X (actually \cancel plus \bcancel), \cancelto{〈value〉}{〈expression〉} draws a diagonal arrow through the 〈expression〉pointing to the 〈value〉 (math-mode only)
\usepackage{csquotes}
\usepackage{afterpage}
\usepackage{parskip} 
\usepackage{float}
\usepackage{enumitem}
%\usepackage{multicol}%Paquete que permite la creación aislada de columnas de texto. De esta forma se reduce la cantidad de páginas en nuestro documento
\usepackage{lipsum} %Paquete que genera texto de relleno al igual que \usepackage{blindtext}, se usa con el comando \lipsum[#-#] los #(hashtags) delimitan el número de páginas que deseamos ocupar del paquete lipsum para usarlos como relleno. 

% Definir el comando personalizado para la nota
\newcommand{\mynote}[1]{%
\begin{tikzpicture}[baseline=-0.75ex]
\node[draw, circle, fill=Apple Green!60, inner sep=2pt] (note) {\textbf{Nota:}};
\end{tikzpicture}%
\ \textit{#1}%
}

% Definición de comandos para variables con barra
\newcommand{\xbar}{\overline{x}}
\newcommand{\ybar}{\overline{y}}

\newenvironment{Figura}
{\par\medskip\noindent\minipage{\linewidth}}
{\endminipage\par\medskip}
\usepackage{caption}
\usepackage[
backend=biber,
sorting=none,
url=true
]{biblatex} %bibliografía
\addbibresource{biblio.bib}
% Establecer el espacio entre las entradas de la bibliografía
\setlength{\bibitemsep}{1\baselineskip} % Puedes ajustar el valor según tus preferencias
\usepackage{amsmath,amsthm,amssymb,amsfonts}
\usepackage{pifont} %Permite la compilación del comando \xmark: es el símbolo de la tachita
\usepackage{mathtools} %permite la compilación de símbolos para matrices
\usepackage{empheq} %Paquete que se relaciona con \usepackage[most]{tcolorbox} para la creación de cuadros/cajas de colores para delimitar los resultados matemáticos o líneas de texto usando la línea de comando: \begin{empheq}[box={\mymath[colback=orange(webcolor)!70,drop lifted shadow]}]{equation*} \end{empheq}
\usepackage[most]{tcolorbox} %Paquete que se relaciona con \usepackage{empheq} para la creación de cuadros/cajas de colores para delimitar los resultados matemáticos o líneas de texto
\renewcommand{\qedsymbol}{$\blacksquare$}
\usetikzlibrary{positioning,decorations.pathreplacing} %Paquete que permite la compilación de llaves para cuadros sinópticos (brace diagram) 
\usepackage{schemata} %The schemata package is designed just for these brace diagrams
\usepackage{booktabs}

\newcommand\AB[2]{\schema{\schemabox{#1}}{\schemabox{#2}}} %comando o paquete necesario para crear cuadros sinópticos

\newtcbox{\mymath}[1][]{%
nobeforeafter, math upper, tcbox raise base,
enhanced, colframe=blue!30!black,
colback=blue!30, boxrule=1pt,
#1} %Comando importante para encasillar los resultados matemáticos en cajas de diferentes colores, se relaciona con los paquetes \usepackage{empheq} y \usepackage[most]{tcolorbox} 

\newenvironment{sysmatrix}[1]
{\left(\begin{array}{@{}#1@{}}}
{\end{array}\right)}
\newcommand{\ro}[1]{%
\xrightarrow{\mathmakebox[\rowidth]{#1}}%
}
\newlength{\rowidth}% row operation width
\AtBeginDocument{\setlength{\rowidth}{3em}}  %Comando importante para laproducción de lineas de operaciones entre matrices, método de Gauss o Gauss Jordan se relaciona con \newenvironment{sysmatrix}

%formato para cambiar el horario a español
\usepackage[spanish]{datetime2}
\DTMsetdatestyle{spanish}

\renewcommand{\today}{\DTMdisplaydate{\the\year}{\the\month}{\the\day }{-1}}
%%%%%%%%%%%%%%%%%%%%%%%%%%%%%%%%%%%%%%%%%

\usepackage{datetime}
\newcommand{\mycurrenttime}{\xxivtime}

%%%%%%%%%%%%%%%%%%%%%%%%%%%%%%%%%%%%%%%%%%%%%%%%%%%%%%%%%%%%%%%%%%%%%%%%%%%%%%%%%%%%%%%%%%%%%%%%%%%%%%%%%%%%%%%%%%%%%%%%%%%%%%%%%%%%%%%%%%%%%%%%%


%%%%%%%%%%%%%Caja de color para el título%%%%%%%%%%%%%%%%%%%%%%%%%%%%%%%%%%%%%%%%%%%%%%%%%%%%%%

\definecolor{myframecolor}{RGB}{1, 45, 75} % Prussian Blue
\definecolor{myboxcolor}{RGB}{0, 107, 92} % Apple Green

\newtcolorbox{mybox}{
enhanced,
colback=myboxcolor!7, % Color de fondo del cuadro
colframe=myframecolor, % Color del marco del cuadro
arc=0pt,
boxrule=1pt,
borderline west={2mm}{-10mm}{myframecolor}, % Borde en el lado izquierdo
sharp corners=southwest,
width=\linewidth,
before=\par\vspace{\bigskipamount}, % Espacio antes del cuadro
after=\par\vspace{\bigskipamount} % Espacio después del cuadro
}

%%%%%%%%%%%%%%%%%%%%%%%%%%%%%%%%%%%%%%%%%%%%%%%%%%%%%%%%%Nueva caja%%%%%%%%%%%%%%%%%%%%%%%%%%%%%%%%%%%%%%%%%%%%%%%%%%%%%%%%%%%%%%%%%

\newtcolorbox{mybox2}{
enhanced,
colback=Apple Green!7, % Color de fondo del cuadro
colframe=Apple Green, % Color del marco del cuadro
arc=0pt,
boxrule=1pt,
borderline west={2mm}{-10mm}{Apple Green}, % Borde en el lado izquierdo
sharp corners=southwest,
width=\linewidth,
before=\par\vspace{\bigskipamount}, % Espacio antes del cuadro
after=\par\vspace{\bigskipamount} % Espacio después del cuadro
}


\newcommand{\euler}{e} %Comando para producir letra e de euler
\newenvironment{remark}{\par\vfill\footnotesize % Vertical white space above the remark and smaller font size
\begin{list}{}{
	\leftmargin=80pt % Indentation on the left
	\rightmargin=60pt}\item\ignorespaces % Indentation on the right
\makebox[-2.5pt]{\begin{tikzpicture}[overlay]
		\node[draw=Horizon!60,line width=2.5pt,circle,fill=Horizon!25,font=\sffamily\bfseries,inner sep=4pt,outer sep=0pt] at (-19pt,5pt){\textcolor{Horizon}{Nota}};\end{tikzpicture}} % Blue Nota in a circle
\advance\baselineskip -1pt}{\end{list}\vskip5pt} % Tighter line spacing and white space after remark
\usepackage{graphicx}

\usepackage{titling}

\input{colores.tex}

%%Producir signos de cita textual``Asignatura''

\newcommand{\xmark}{\ding{55}} %%%comando que genera la tachita

\newenvironment{MyColorPar}[1]{%
\leavevmode\color{#1}\ignorespaces%
}{%
}%

%%%%%%%%%%%%%%%%%%% Título de la tarea, Nombre de alumno

\title{Clase \# 41 - Estado Sólido}
\author{\emph{JA} $\boldsymbol{\mid}$}
\date{\today}

%%%%%%%%%%%%%%%%%%% Datos de la Materia

\usepackage{fancyhdr}
\fancypagestyle{plain}{%  the preset of fancyhdr 
\fancyhf{} % clear all header and footer fields
\fancyfoot[R]{
\includegraphics[width=2cm]{LogoFCFMBUAP (1).png}}
\fancyfoot[L]{{\bfseries{\thedate{}}} a las {\bfseries{\mycurrenttime{}}} horas{} (GMT-6, H. Puebla de Zaragoza, Pue)}
\fancyhead[L]{Estado sólido}
\fancyhead[R]{\theauthor}
}
\makeatletter
\def\@maketitle{%
\newpage
\null
\vskip 1em%
\begin{center}%
\let \footnote \thanks
{\LARGE \@title \par}%
\vskip 1em%
%{\large \@date}%
\end{center}%
\par
\vskip 1em}
\makeatother

\usepackage{lipsum}  
\usepackage{cmbright}

%%%%%%%%%%%%%%%%%%%%%%%%%%%%%%%%%%%%%%%%%%%%%%%%%%%%%%%%%%%%%%%%%%%%%%%%%%%%%%%%%%%%%%%%%%%%%%%%%%%%%%%%%%%%%%%%%%%%%%%%%%%%%%%%%%%%%%%%%%%%%%%%%%%%%%%%%%%%%%%%%%%%%%%%%%%%%%%%%%%%%%%%%%%%%%%%%%%%%%

% Por si el color se llama "Apple Green" con espacio en colores.tex:
\colorlet{AppleGreen}{Apple Green}

% Caja AppleGreen con tono ajustable
\newenvironment{AGbox}[1][7]% 7 = tono por defecto (más grande => más claro)
{%
  \begin{tcolorbox}[
    enhanced, breakable,
    colframe=AppleGreen,
    colback=AppleGreen!#1,
    boxrule=1pt,
    borderline west={2mm}{-10mm}{AppleGreen},
    sharp corners=southwest,
    left=14pt,right=14pt,top=10pt,bottom=10pt,
     boxsep=6pt,
  ]%
}{%
  \end{tcolorbox}%
}



%%%%%%%%%%%%%%%%%%%%%%%%%%%%%%%%%%%%%%%%%%%%%%%%%%%%%%%%%%%%%%%%%%%%%%%%%%%%%%%%%%%%%%%%%%%%%%%%%%%%%%%%%%%%%%%%%%%%%%%%%%%%%%%%%%%%%%%%%%%%%%%%%%%%%%%%%%%%%%%%%%%%%%%%%%%%%%%%%%%%%%%%%%%%%%%%%%%%%%

\begin{document}



\maketitle


%%%%%%%%%%%%%%%%%%% Datos del profesor, campus y aula

\begin{mybox}
\noindent\begin{tabular}{@{}ll}
	{\bfseries{Profesor}} & Dra. Ana Lilia González Ronquillo\\
	\textcolor{prussianblue}{\bfseries{Campus}} / \textcolor{trueblue}{\bfseries{Facultad}}  & \textcolor{prussianblue}{\bfseries{C.U. BUAP}} /  \textcolor{trueblue}{\bfseries{FCFM}} \\
	\textcolor{prussianblue}{\bfseries{Edificio}} /  \textcolor{trueblue}{\bfseries{Salón}}    & 	\textcolor{prussianblue}{\bfseries{FM9}} /  \textcolor{trueblue}{\bfseries{301}}   
\end{tabular} 
\end{mybox} \vspace{1cm}

\tableofcontents

\vspace{2cm}


\section{Clase de hoy}


Tenemos:


\section{\textcolor{orange(colorwheel)}{Dependencia de la temperatura}}




La integral es del tipo de intefrales de Fermi-Dirac:


\[  F_{j}(y_{0}) = \int\limits_{0}^{\infty}  \frac{ y^{j}  }{ e^{ (y-y_{0}) +1 }  } dy   \]

Una solució aproximada converge cuando $y_{0}>>1$


EN el caso de los metales:

\[  K_{B}T \approx 0\text{.}01\mu \]


Nos pidió que avanzaramos en el resultado de:


\[   n \approx \frac{1}{3\pi^{2}} \left(  \frac{2m}{\hbar^{2}}  \right)^{\frac{3}{2}} \mu^{\frac{3}{2}} \left[   1+ \frac{ \pi^{2}  }{  8 \left(  \frac{ \mu }{ K_{B}T }  \right)^{2} }    \right]   \]

El número de partículas por unidad de volumen no va cambiar.


Por otro lado, sabemos que el número de electrones ñor unidad e volumen no depende de la temperarura, es decir $n$ es constante. Ya teníamos una expresión para $n$:



\[   n = \frac{  k^{3}_{f} }{3\pi^{2}} = \left(  \frac{2m}{\hbar^{2}}\epsilon_{f}  \right)^{\frac{3}{2}} \frac{1}{3\pi^{2}} \]


Las ecuaciones anteiores deben de ser iguales, tenemos así:


\[   \left(   \frac{2m}{\hbar^{2}}  \right)^{\frac{3}{2}} \ldots \, \text{falta contenido, checar notas}     \]


\textbf{Resultado relevante}


Tenemos:


\[ \epsilon^{\frac{3}{2}}  \approx \mu^{\frac{3}{2}}  \left[  1 + \frac{ \pi^{2}  }{  8 \left(   \frac{\mu}{  K_{B}T } \right)^{2} }    \right] \]




O bien, despejamos el \textbf{potencial químico} y se tiene: 


\[  \mu \approx \epsilon_{f}  \left[  1 + \frac{ \pi^{2}  }{ 8 \left( \frac{\mu}{ K_{B}T  }   \right) }  \right]  \]

¿Qué implica la convergencia de la serie?


Recordandoq ue $K_{B}T \approx 0\text{.}01\mu$, podemos usar $(1+x)^{r}\approx 1+rx$ con $x<<1$


Tenemos:



\[ \mu \approx \epsilon_{f} \left[1 - \frac{  (\pi K_{B}T)^{2}  }{12\mu^{2}}    \right]   \]


Hemos dicho que cuando $T$ tiende a cero, el \textbf{potencial químico} tiende al nivel de \textbf{Fermi}, es decir:


\[   \lim\limits_{T \to } \mu = \epsilon_{f} \]

Como aproximación podemos reemplazar $\mu \to \epsilon_{f}$ en el tpermino proporcional a $T^{2}$,  así finalmente tenemos:


\[ \mu \approx \epsilon_{f} \left[    \right]  \]



De las cosas que hemos hecho con Sommerfield, tenemos nosotros:


NOs interesa ahora es una temperatura diferente de cero, tenemos que resolver esta integral (SI es cero, es una función de paso).

\[ \boxed{ \text{ a $T\neq 0^{\circ}K$ } } \]

\[  U/V = \int\limits_{0}^{\infty} \epsilon g(\epsilon) f(\epsilon)d\epsilon = \frac{1}{2}  \]

Tendremos que usar la serie de $F_{j}(y_{0} = \frac{}{})$


Hay que considerar lo que vimos.

Donde nos conviene cambiar $\mu$, por el nivel de Fermi. 

De manera general vamos a tener en un cristal:


\[ C_{cristal} = AT + BT^{3} \]

LOs semicondutores son tetravalantes, quiero decir que los enclazamos con otros cuatro atomos.

LOs enlaces, tetraedro. 

ENlaces localizados, la distriucuión de carga, en el metal es una estructura periodpica, huecos, electrones y huecos, que se muevem.

\textcolor{Cinnabar}{\textbf{Próximos exámenes}}:

\textbf{Examen de recuperación}: $\boxed{\text{3 de diciembre}}$ (será el primer parcial dado que hay mayor cantidad de no aprobados)

\textbf{Examen de teoría de electrones y teoría de bandas (últimos dos temas)}:  $\boxed{\text{1 de diciembre}}$. 

Las calificaciones finales las mandará el día domingo 7 de diciembre, podemos ir a su cubículo en el \underline{Instituto} por alguna duda. 

\section{Nuevo tema: Teoría de bandas}

Este tema es muy importante. La clasificación de los materiales, nos vamos a enfocar en la teoría de bandas, que explica teoría de bandas y de semiconductores, es una imagen del KIttel, es un diagra de las diapositivas que aún va a subir.


Investigar qué es un \textbf{semimetal}. Los materiales 

Recorar la escrutura del diamante dormada pror carnono, los benos los bamos a llamar, aislasntes. 

\begin{figure}[H]
    \centering
    \includegraphics[width=0.75\textwidth]{1.png}
    \caption{Aquí va la descripción de tu figura.}
    \label{fig:mi_figura_png}
\end{figure}


Existen dos tipos de semiconductores, los \textbf{íntrinsecos} y los \textbf{extrinsecos}.

¿Qué son las bandas de energía?


SOn conjuntos de niveles de energía. 

Vamos a ver un modelo de por qué aparecen las bandas en un material periódico.

Van a aparecer niveles que son permitidos, cuales no y así. 



En el caso de los aislante resulta que tenemos muchas bandas permitidas, la última banda en un material aislante está completamente llena, la siguiente banda está vacía.

Necesitamos calentarlo lo suficiente para que puedan hacer el salto, 

Resulta que en los materiales aislantes. 


En los metales, la última banda estásn semillenos. Los metales sufren \textbf{Transiciones intrabanda}.


NO es necesario calentarlos tanto para llegar a la última banda. 


El otro caso es el del caso de los \textbf{semiconductores}. En ese material tenemos bandas permitidas y prohibidas. 


\begin{figure}[H]
    \centering
    \includegraphics[width=0.75\textwidth]{2.png}
    \caption{Aquí va la descripción de tu figura.}
    \label{fig:mi_figura_png}
\end{figure}

\begin{figure}[H]
    \centering
    \includegraphics[width=0.75\textwidth]{3.png}
    \caption{Aquí va la descripción de tu figura.}
    \label{fig:mi_figura_png}
\end{figure}


Nosotros estamos graficando esto:

\[ \text{Relación de dispersión}  \]

\begin{figure}[H]
    \centering
    \includegraphics[width=0.75\textwidth]{4.png}
    \caption{Aquí va la descripción de tu figura.}
    \label{fig:mi_figura_png}
\end{figure}

\begin{figure}[H]
    \centering
    \includegraphics[width=0.75\textwidth]{5.png}
    \caption{Aquí va la descripción de tu figura.}
    \label{fig:mi_figura_png}
\end{figure}

\begin{figure}[H]
    \centering
    \includegraphics[width=0.75\textwidth]{6.png}
    \caption{Aquí va la descripción de tu figura.}
    \label{fig:mi_figura_png}
\end{figure}

\begin{figure}[H]
    \centering
    \includegraphics[width=0.75\textwidth]{7.png}
    \caption{Aquí va la descripción de tu figura.}
    \label{fig:mi_figura_png}
\end{figure}



Tenemos (recordando espacio recíproco) :


\[  v(\vec{r}) = \sum\limits_{\vec{G}} V_{\vec{G}} e^{i \vec{G} \cdot \vec{r} }  \]




\[  \vec{G} = m_{1}\vec{b}_{1} + m_{2}\vec{b}_{2} + m_{3}\vec{b}_{3}  \]



Tenemos:


\[  \sum\limits_{\vec{G}} \to \sum\limits_{m_{1}= - \infty}\sum\limits_{m_{2}= - \infty}\sum\limits_{m_{3}= - \infty}       \]


\begin{figure}[H]
    \centering
    \includegraphics[width=0.75\textwidth]{8.png}
    \caption{Aquí va la descripción de tu figura.}
    \label{fig:mi_figura_png}
\end{figure}




\end{document}
